%Chargement des paquets
\usepackage{amsmath}
\usepackage{amsfonts}
\usepackage{amssymb}
\usepackage{enumerate}
\usepackage{mathtools}
\usepackage{bbm}
\usepackage{xparse, etoolbox}
\usepackage{enumerate}
\usepackage{enumitem}
\usepackage{mathabx}
\usepackage{babel}[french]

%environment
\usepackage{amsthm}
\usepackage{keytheorems}
\usepackage{hyperref} 

%Interface théorème
\newkeytheorem{lem}[
	name = Lemme
]

\newkeytheorem{prop}[
	name = Proposition
]

\newkeytheorem{thm}[
	name = Théorème
]

\newkeytheorem{coro}[
	name = Corollaire
]

\renewcommand*{\proofname}{Démonstration}

\theoremstyle{definition}
\newtheorem{defn}{Définition}

\theoremstyle{definition}
\newtheorem*{nota}{Notation}

\theoremstyle{definition}
\newtheorem*{limi}{Liminaire}

\theoremstyle{remark}
\newtheorem{rmq}{Remarque}

\theoremstyle{plain}
\newtheorem*{fait}{Fait}


%Ensembles de nombres
\newcommand{\N} {\mathbb{N}}
\newcommand{\Ne}{\N^\ast}
\newcommand{\Z} {\mathbb{Z}}
\newcommand{\Q} {\mathbb{Q}}
\newcommand{\R} {\mathbb{R}}
\newcommand{\Rb}{\overline{\mathbb{R}}}
\newcommand{\Rp}{\R_+}
\newcommand{\Rm}{\R_-}
\newcommand{\K} {\mathbb{K}}
\newcommand{\Ket} {\mathbb{K^{\ast}}}
\newcommand{\Cx}{\mathbb{C}}

%Ensembles, tribus
\newcommand{\A}{\mathbb{A}}
\renewcommand{\P}{\mathcal{P}}
\newcommand{\T}{\mathcal{T}}
\newcommand{\F}{\mathcal{F}}
\newcommand{\C} {\mathcal{C}}
\newcommand{\ouv}{\mathcal{O}}
\newcommand{\B}{\mathcal{B}}
\newcommand{\M}{\mathcal{M}}
\newcommand{\etag}{\mathcal{E}}
\newcommand{\etagOp}{\etag^{+}(\Omega)}
\newcommand{\etagO}{\etag(\Omega)}
%%Lp avant quotientage
\newcommand{\Lic}{\mathcal{L}^1}
\newcommand{\Lpc}{\mathcal{L}^p}
\newcommand{\Linfc}{\mathcal{L}^{\infty}}
\newcommand{\Lifull}{\Lic(\Omega, \B, m)}
\newcommand{\Lpfull}{\Lpc(\Omega, \B, m)}
\newcommand{\Linffull}{\Linfc(\Omega, \B, m)}
%%Lp après quotientage
\newcommand{\Li}{L^1}
\newcommand{\Ld}{L^2}
\newcommand{\Lp}{L^p}
\newcommand{\Linf}{L^{\infty}}
\newcommand{\Listd}{\Li(\Omega, m)} 
\newcommand{\Ldstd}{\Ld(\Omega, m)} 
\newcommand{\Lpstd}{\Lp(\Omega, m)} 

%Quelques opérateurs
\DeclareMathOperator{\codim}{codim}
\DeclareMathOperator{\rg}{rg}
\DeclareMathOperator{\tr}{tr}
\DeclareMathOperator{\vect}{Vect}
\DeclareMathOperator{\ima}{Im}
\DeclareMathOperator{\id}{Id}
\DeclareMathOperator{\sign}{sign}
\DeclareMathOperator{\ext}{ext}
\newcommand{\ind}{\mathbbm{1}}
\newcommand*{\spr}[2]{\left(\,#1\,\middle|\,#2\,\right)}
\newcommand{\interior}[1]{%
  {\kern0pt#1}^{\mathrm{o}}%
}
\DeclareMathOperator{\card}{card}
\DeclareMathOperator{\ppcm}{ppcm}

%Commandes de pure flemme
\newcommand{\veps}{\varepsilon}
\newcommand{\intab}{\int_{a}^{b}}
\newcommand{\intR}{\int_{-\infty}^{\infty}}
\newcommand{\intinf}{\int_{- \infty}^{+ \infty}}
\newcommand{\equi}{\Leftrightarrow}
\newcommand{\Kev}{$\K$-espace vectoriel }
\newcommand{\Kevs}{$\K$-espaces vectoriels }
\newcommand{\Rev}{$\R$-espace vectoriel }
\newcommand{\Revs}{$\R$-espaces vectoriels }
\newcommand{\Cev}{$\Cx$-espace vectoriel }
\newcommand{\Cevs}{$\Cx$-espaces vectoriels }
\newcommand{\evn}{(E, \norm{.})}
\newcommand{\interf}[2]{[#1,#2]}
\newcommand{\limb}{\lim\limits_}
\newcommand{\lebn}{\lambda_n}
\newcommand{\pp}{\text{~presque partout}}
\newcommand{\derivp}[2]{\frac{\partial #1}{\partial #2}}
\newcommand{\famiu}{(u_i)_{i \in I}}
\newcommand{\sev}{\text{~s.e.v.}}
\newcommand{\Mnp}{M_{n,p}(\K)}
\newcommand{\Mn}{M_{n}(\K)}
\newcommand{\tq}{\text{~t.q.~}}
\newcommand{\mq}{~montrer~que~}
\newcommand{\et}{\quad \text{et} \quad}
\newcommand{\ou}{\text{~ou~}}
\newcommand{\Kx}{\K[X]}


%Commandes de gars chelou
\renewcommand{\subset}{\subseteq}

%Commande restriction, ty duckduckgo
\def\restr#1#2{\mathchoice
              {\setbox1\hbox{${\displaystyle #1}_{\scriptstyle #2}$}
              \restraux{#1}{#2}}
              {\setbox1\hbox{${\textstyle #1}_{\scriptstyle #2}$}
              \restraux{#1}{#2}}
              {\setbox1\hbox{${\scriptstyle #1}_{\scriptscriptstyle #2}$}
              \restraux{#1}{#2}}
              {\setbox1\hbox{${\scriptscriptstyle #1}_{\scriptscriptstyle #2}$}
              \restraux{#1}{#2}}}
\def\restraux#1#2{{#1\,\smash{\vrule height .8\ht1 depth .85\dp1}}_{\,#2}} 

%Quelques normes
\DeclarePairedDelimiter\abs{\lvert}{\rvert}
\DeclarePairedDelimiter\labs{\bigl\lvert}{\bigr\rvert}
\DeclarePairedDelimiter\norm{\lVert}{\rVert}
\DeclarePairedDelimiterXPP\normi[1]{}\lVert\rVert{_1}{\ifblank{#1}{\:\cdot\:}{#1}}
\DeclarePairedDelimiterXPP\normd[1]{}\lVert\rVert{_2}{\ifblank{#1}{\:\cdot\:}{#1}}
\DeclarePairedDelimiterXPP\normp[1]{}\lVert\rVert{_p}{\ifblank{#1}{\:\cdot\:}{#1}}
\DeclarePairedDelimiterXPP\normq[1]{}\lVert\rVert{_q}{\ifblank{#1}{\:\cdot\:}{#1}}
\DeclarePairedDelimiterXPP\norminf[1]{}\lVert\rVert{_\infty}{\ifblank{#1}{\:\cdot\:}{#1}}

%Bracket en angle
\DeclarePairedDelimiter\angl{\langle}{\rangle}

%Set, parenthèses
\newcommand{\set}[1]{\{ #1 \}}
\newcommand{\bigset}[1]{\bigl\{ #1 \bigr\}}
\newcommand{\Bigset}[1]{\Bigl\{ #1 \Bigr\}}
\newcommand{\bigpar}[1]{\bigl( #1 \bigr)}
\newcommand{\Bigpar}[1]{\Bigl( #1 \Bigr)}

%Opitmisation des opérateurs de comparaison
\DeclarePairedDelimiter\tio{o (}{)}
\DeclarePairedDelimiter\gro{O (}{)}

\newcommand{\Mnr}{M_n(\R)}
\newcommand{\So}{\text{SO}(n,\R)}
\newcommand{\Sl}{\text{SL}(n,\R)}

\pagestyle{fancy}

\fancyhead[C]{}
\fancyhead[R]{}

\begin{document}

\underline{Source :} Gourdon - Analyse - page 341

\underline{Leçons :}
\begin{itemize}
\item 
\end{itemize}

\begin{lem}[Inégalité de Cauchy-Schwarz, cas euclidien]
Soit $E$ un espace euclidien (de fimension finie $m$), dont le produit scalaire est dénoté par $\angl{\cdot, \cdot}$.
Pour $x, y \in E$ on a :
\[ \abs{\angl{x,y}} \leq \norm{x} \norm{y} , \]
où $\norm{\cdot}$ est la norme issue du produit scalaire.
\end{lem}

\begin{proof}
Soient $x, y \in E$, avec $y$ non nul et $t \in \R$, on a :
\begin{align*}
\angl{x+ty, x+ty} & = \norm{x+ty} \\
& = \norm{x}^2 + t^2 \norm{y}^2 + 2t \angl{x,y}
& = P(t)
\end{align*}
Il s'agit donc d'un polynôme de degré 2 en $t$ qui, par définition, est positif ou nul pour tout $t$. Son discriminant est donc négatif ou nul,
soit :
\[ 4 \angl{x,y} - 4 \norm{x}^2 \norm{y}^2 \leq 0 , \]
ce qui donne directement le résultat. Dans le cas où $y=0$, l'inégalité est triviale.
\end{proof}

\begin{lem}[Continuité de la norme]
Soit $E$ un \Kev de dimension finie et $N$ une norme sur $E$. Alors $N$ est une application continue de $E$ dans $\R$.
\end{lem}

\begin{thm}[Passer à l'oral]
Soit $f : U \subset \R^n \to F, n \geq 1$, une application, avec $U$ un ouvert de $\R^n$ et $F$ un e.v.n.
Si toutes les dérivées partielles de $f$ sur $U$ existent et si elles sont continues en un point $a$ de $U$, 
alors $f$ est différentiable en $a$ et on a :
\[ df_a = \sum_{i=1}^n \dfrac{\partial f}{\partial x_i} (a) dx_i , \]
où $(dx_1, \dots, dx_n)$ est la base duale (canonique) de $\R^n$.
\end{thm}

\begin{proof}
$\R^n$ étant de dimension fini, on peut le normer de façon quelconque, on choisit donc la norme $\norm{\cdot}$ définie par :
\[ \forall x \in \R^n, \quad \norm{x} = \sum_{i=1}^n \abs{x_i} . \]
On va étudier l'application $g$ définie par :
\[ \forall x \in U, \quad g(x) = f(x) - \sum_{i=1}^n \dfrac{\partial f}{\partial x_i} (a) x_i , \]
en particulier, on cherche à montrer que :
\[ g(x) - g(a) = f(x) - f(a) - \sum_{i=1}^n (x_i-a_i) \dfrac{\partial f}{\partial x_i} (a) = o(\norm{x-a}) , \]
au voisinage de $a$ (i.e. pour $\norm{x-a} \to 0$).
Soit $\veps>0$, par hypothèse de continuité sur les dérivées partielles on peut écrire :
\[ \exists \alpha >0, \forall x \in U, \norm{x-a} < \alpha, \forall i \in \set{1, \dots, n}, \quad 
	\norm*{\dfrac{\partial g}{\partial x_i}(x)} = \norm*{\dfrac{\partial f}{\partial x_i}(x) - \dfrac{\partial f}{\partial x_i}(a)} < \veps . \]
De plus, quitte à réduire $\alpha$, on peut toujours supposer que $B(a, \alpha) \subset U$.
Prenons maintenant $x \in B(a, \alpha)$, et considérons les points :
\[ y_0 = (a_1, \dots, a_n) \et \forall k \in \set{1, \dots, n} y_k = (x_1, \dots, x_k, a_{k+1}, \dots, a_n) , \]
de sorte que $y_0 = a$ et $y_n=x$. Pour $k \in \set{1, \dots, n}$, on définit la fonction $g_k$ par :
\[ \begin{array}{c c c c}
g_k : & (a_k, x_k) & \to & F \\
& t & \mapsto & g(x_1, \dots, x_{k-1}, t, a_{k+1}, a_n) , 
\end{array} \]
où $(a_k, x_k)$ est l'intervalle de bornes $a_k$ et $x_k$. Cette fonction vérifie :
\[ g'_k(t) = \dfrac{\partial g}{\partial x_k}(x_1, \dots, x_{k-1}, t, a_{k+1}, a_n) , \]
donc clairement, $\norm{g'_k(t)} < \veps$ sur $(a_k, x_k)$. On en déduit que :
\[ \norm{g_k(x_k) - g_k(a_k)} \leq \veps \abs{x_k-a_k} , \]
or $g_k(x_k) = g(y_k)$ et $g_k(a_k) = g(y_{k-1})$ donc :
\[ \norm{g(y_k) - g(y_{k-1})} \leq \veps \abs{x_k-a_k} , \]
et par suite :
\[ \norm{g(x) - g(a)} = \norm*{\sum_{k=1}^n (g(y_k) - g(y_{k-1}))} \leq sum_{k=1}^n \norm{g(y_k) - g(y_{k-1})}
	\leq \veps sum_{k=1}^n \abs{x_k-a_k} = \veps \norm{x-a} . \]
Et ainsi, on a montré que $g(x)-g(a) \in o(\norm{x-a})$ au voisinage de $a$.
\end{proof}

\begin{lem}[Régularité du déterminant]
L'application $\det$ définie sur $\Mnr, n \geq 1$ et à valeur dans $\R$ est de classe $\C^{\infty}$.
\end{lem}

\begin{proof}
L'espace vectoriel $\Mnr$ est de dimension finie, on peut donc le munir d'une norme quelconque, par exemple la norme 
$\norminf{\cdot}$. Le déterminant est une fonction polynômiale des coefficients de la matrice, donc une application $\C^{\infty}$.
\end{proof}

\begin{lem}[Différentielle du déterminant]
L'application $\det$ admet une différentielle en tout point $M \in \Mnr$ que l'on note $D \det_M$ et on a : 
\[ D \det_M : H \mapsto \tr(C^T H) , \]
où $C=(C_{i,j)}$ est la matrice des cofacteurs de la matrice $M$.
\end{thm}

\begin{proof}
La différentiabilité de l'application déterminant est conséquence dans la continuité de toutes ses dérivées partielles 
(le déterminant d'une matrice est un polynôme en les coefficients de cette dernière).
Soit $M=(m_{i,j}) \in \Mnr$, on va calculer les dérivées partielles de l'application $\det$ au point $M$.
Pour cela, on note $A=(a_{i,j})$ la forme générale des éléments de $\Mnr$, 
et $(E_{i,j})_{1 \leq i,j \leq n}$ la base canonique de $\Mnr$.
On va calculer :
\[ \dfrac{\partial \det}{\partial a{i,j}}(M) , \]
pour chaque couple d'indices $(i,j)$. On a :
\[ \dfrac{\partial \det}{\partial a{i,j}}(M) = \lim_{\substack{t \to 0 \\ t \neq 0}} \dfrac{ \det(M + t E_{i,j}) - \det{M}}{t} = C_{i,j} . \]
Donc, pour $H = (h_{i,j}) \in \Mnr$ on a :
\[ D \det_M = \sum_{i,j} \dfrac{\partial \det}{\partial a{i,j}}(M) h_{i,j} = \sum_{i,j} h_{i,j} C_{i,j} = \tr(C^TH) . \]
\end{proof}

\begin{thm}
Soit $n \in \Ne$, on munit $\Mn$ de la norme :
\[ \forall M=(m_{i,j}) \in \Mn, \quad \norm{M} = (\sum_{i,j} m_{i,j}^2)^{1/2} . \]
Le groupe des matrices orthogonales directes $\So$ est l'ensemble des éléments de $\Sl$ de norme minimale.
\end{thm}

\begin{proof}
On va déjà montrer que l'application définie est bien une norme sur $\Mn$.
\begin{enumerate}[label=(\roman*)]
\item Soit $M \in \Mn$ telle que $\sum_{i,j} m_{i,j}^2 = 0$, il est alors clair que $M$ est la matrice nulle.
\item Soit $M \in \Mn$ et $\lambda \in \R$, par calcul direct on a bien :
\[ \norm{\lambda M} = (\sum_{i,j} \lambda^2 m_{i,j}^2)^{1/2} = \abs{\lambda} \norm{M}. \]
\item La sous-additivité de l'application $\norm{.}$ est une conséquence directe de l'inégalité de Cauchy-Schwarz
(en identifiant une matrice de $\Mn$ à un vecteur de $\R^{n^2}$).
\end{enumerate}
Ceci étant fait, on remarque que $\norm{M}^2 = \tr(M^TM)$, ainsi le résultat du théorème est un problème de minimisation
de la fonction $f : M \mapsto \tr(M^TM)$ sous la contrainte $g(M) = 0$ avec $g: M \mapsto \det(M)-1$.
Il est clair que $f$ et $g$ sont des applications continues. En particulier, l'ensemble :
\[ \Sln =g^{-1}(\set{0}) , \]
est un fermé de $\Mnr$. Notons $d = \inf \set{\norm{M} ; M \in \Sln}$, alors l'ensemble $\Delta = \Sln \cap B_f(0,d+1)$ est un fermé
borné de $\Mnr$ (ici $B_f(0,d+1)$ est une boule fermée dans $\Mnr$) donc un compact.
Par suite, la fonction $f$ atteint sa borne inférieure sur ce compact. 
Notons $A \in \Sln$ la matrice minimisant $f$ sur $\Sln$.
La différentielle de la fonction $g$ en $A$ est l'application $dg_A$ définie par :
\[ \forall H \in \Mnr, \quad dg_A = \tr(C^TH) , \]
où $C$ est la comatrice de $A$.
\end{proof}
 
\end{document}